\documentclass[11pt,a4paper,twoside]{book}
\usepackage{geometry}
\geometry{ %
    paperwidth=170mm, 
    paperheight=220mm, 
    inner=30mm,   % Increased gutter for binding
    outer=25mm,   % Outer margin slightly smaller
    top=20mm,     % More breathing space for headers
    bottom=25mm,  % More space for footers and page numbers
    headheight=14pt,  % Adjusts space for headers
    footskip=12mm     % Space between text and footer
}
\usepackage[greek,english]{babel}
\usepackage[utf8]{inputenc}
\usepackage[T1]{fontenc}
\usepackage{hyperref}
\hypersetup{pageanchor=true}
\usepackage{amsmath, amssymb}
\usepackage{booktabs} % table fanciness
\usepackage{algorithm} % For the algorithm environment
\usepackage{algorithmic} % For the algorithmic pseudocode formatting
\usepackage{mdframed}  % 
\usepackage{amsthm}  %  theorem-style environments
\usepackage{bussproofs} % graphical theorem/axioms for natural deduction
\usepackage{tikz} % tree
\usepackage{libertine} % override lmodern?  libertine
\usepackage{pifont} %
\renewcommand{\familydefault}{\sfdefault}
\usepackage{enumerate}
\usepackage{listings} % code
\usepackage{tabularx} % tables
\usetikzlibrary{positioning,arrows.meta} % nice arrows?
\usepackage{verbatim} % verbatim smaller
\usepackage{fancyhdr} %
\usepackage{tocloft} %
\usepackage{imakeidx} % index
\usepackage{idxlayout} % index
\usepackage{titlesec} % headlines
\usepackage{endnotes} %
\usepackage{marginnote}
\usepackage{xcolor} %
\usepackage{framed}
\usepackage{graphicx} % 
\usepackage{pdfx} %
\usepackage{xr}
\usepackage{qtree} %
\usetikzlibrary{positioning, shapes.multipart, arrows.meta, calc} %
\usepackage{microtype} % fix e.g. crappy line breaks
\usepackage{caption}  %
\usepackage{amsthm}     % for the 'definition' environment
\usepackage{morewrites}
\usepackage{caption}
	\captionsetup{font=small}
\usepackage{afterpage}
\usepackage{listings}
\usepackage[scaled]{inconsolata}

\bibliographystyle{apalike}
\pagestyle{plain} %

\newtheorem{definition}{Definition}
\newtheorem{example}{Example}
\newtheorem{theorem}{Theorem}
\newtheorem{abstract}{Abstract}

% You do not need all of this! You probably can get a long way without most of the above ..
% And there are overflows .. that you can fix with a better preamble!

\title{The Ferrite Language}
%\author{}
%\date{\today}

\begin{document}
\maketitle


\tableofcontents
\bigskip


\section{The Ferrite Language: Formal Specification}

\subsection{Overview}

Ferrite is a statically typed, expression-based language with explicit ownership
and borrowing semantics inspired by Rust. Its design philosophy centers on
\emph{compile-time memory safety} without requiring garbage collection or manual
memory management. The semantics are defined by a combination of typing rules and
small-step operational semantics, where the type system enforces memory safety
through precise ownership tracking.

The language features an affine type system, meaning values can be used at most
once unless they implement the \texttt{Copy} trait. This restriction, combined
with explicit borrowing, prevents common memory safety violations such as
use-after-free, double-free, and data races at compile time.


\subsection{Judgement Forms}

The formal semantics of Ferrite are expressed through two primary judgement forms:

\paragraph{Typing Judgements}
\[
\Gamma ; \Sigma \vdash e : \tau
\]
This judgement reads: ``Under type environment $\Gamma$ and ownership environment
$\Sigma$, expression $e$ has type $\tau$.'' The separation of type and ownership
information is crucial---$\Gamma$ tracks \emph{what} types variables have, while
$\Sigma$ tracks \emph{how} they can be used (owned, moved, or borrowed).

\paragraph{Evaluation Judgements}
\[
\langle e, \Sigma \rangle \rightarrow \langle e', \Sigma' \rangle
\]
This judgement specifies that expression $e$ in ownership state $\Sigma$ evaluates
in one step to expression $e'$ in ownership state $\Sigma'$. Notably, the
ownership state can change during evaluation (e.g., when a value is moved or
borrowed), while types remain invariant (Preservation theorem).


\subsection{Types}

\[
\tau ::= i32 \mid bool \mid unit \mid S \mid \&\tau \mid \&_{mut}\tau
\]

Where $S$ ranges over structure types (user-defined structs).

\paragraph{Type Descriptions}
\begin{itemize}
  \item $i32$: 32-bit signed integer (primitive, \texttt{Copy})
  \item $bool$: Boolean value (primitive, \texttt{Copy})
  \item $unit$: Empty type, analogous to \texttt{void} in C or \texttt{()} in Rust
  \item $S$: Struct type with named fields
  \item $\&\tau$: Immutable reference to a value of type $\tau$ (shared borrow)
  \item $\&_{mut}\tau$: Mutable reference to a value of type $\tau$ (exclusive borrow)
\end{itemize}

Reference types are \emph{non-owning} pointers with lifetime constraints enforced
at compile time. Unlike owned types, references cannot outlive their referents.


\subsection{Environments}

Ferrite maintains two distinct but coordinated environments during type checking:

\paragraph{Type Environment}
\[
\Gamma ::= \emptyset \mid \Gamma, x : \tau
\]
Maps variables to their types. This is the traditional typing context found in
most typed languages. Once bound, a variable's type never changes (though it may
become inaccessible after being moved).

\paragraph{Ownership Environment}
\[
\Sigma ::= \emptyset \mid \Sigma, x : \omega
\]
\[
\omega ::= owned \mid moved \mid borrowed(n)
\]
Tracks the ownership state of each variable:
\begin{itemize}
  \item $owned$: The variable owns its value and can be freely used or moved
  \item $moved$: The value has been transferred elsewhere; the variable is
    inaccessible
  \item $borrowed(n)$: The value has $n$ active immutable borrows; it can be read
    but not moved or mutably borrowed
\end{itemize}

The ownership environment $\Sigma$ changes during evaluation and at scope
boundaries, reflecting the runtime behavior of the ownership system.


\subsection{Values}

\[
v ::= n \mid true \mid false \mid S(v_1,\dots,v_n) \mid \&x \mid \&_{mut}x
\]

Values are fully evaluated expressions:
\begin{itemize}
  \item $n$: Integer literals
  \item $true$, $false$: Boolean literals
  \item $S(v_1,\dots,v_n)$: Struct value with fields initialized to values
  \item $\&x$: Immutable reference to variable $x$
  \item $\&_{mut}x$: Mutable reference to variable $x$
\end{itemize}

Note that references are \emph{syntactic}: they refer to variables in scope,
not arbitrary heap addresses. This simplification enables straightforward
lifetime checking.


\subsection{Typing Rules}

The typing rules define when an expression is well-typed. Each rule has premises
(above the line) that must hold for the conclusion (below the line) to be valid.

\paragraph{Variables}

\[
\frac{
  \Gamma(x) = \tau \quad \Sigma(x) = owned
}{
  \Gamma ; \Sigma \vdash x : \tau
}
\quad \textsc{T-Var}
\]

A variable can be used only if it is currently owned. This rule prevents
use-after-move errors. Note that using a variable in an owned context will
trigger a move during evaluation (see operational semantics).

\paragraph{Copy Types}

\[
\frac{
  \Gamma(x) = \tau \quad \tau \in \mathit{Copy} \quad \Sigma(x) \neq moved
}{
  \Gamma ; \Sigma \vdash x : \tau
}
\quad \textsc{T-Var-Copy}
\]

For types implementing \texttt{Copy} (primitives and references), the variable
can be used multiple times without moving. The ownership state remains unchanged.

\paragraph{Let Binding}

\[
\frac{
  \Gamma ; \Sigma \vdash e_1 : \tau_1 \quad
  \Gamma[x:\tau_1] ; \Sigma[x \mapsto owned] \vdash e_2 : \tau_2
}{
  \Gamma ; \Sigma \vdash (\mathtt{let}\ ((x\ e_1))\ e_2) : \tau_2
}
\quad \textsc{T-Let}
\]

A let-binding evaluates $e_1$ to produce a value, binds it to $x$ with ownership,
then evaluates the body $e_2$ in the extended environment. The type of the entire
expression is the type of the body. If $e_1$ consumes any variables, they are
marked as moved in $\Sigma$ before type-checking $e_2$.

\paragraph{Struct Construction}

\[
\frac{
  \Gamma ; \Sigma \vdash e_i : \tau_i \quad \forall i \quad
  \mathit{fields}(S) = \{f_i : \tau_i\}
}{
  \Gamma ; \Sigma \vdash (S\ e_1\ \dots\ e_n) : S
}
\quad \textsc{T-Struct}
\]

Constructing a struct requires providing expressions for each field. Each
argument $e_i$ is evaluated and its value is \emph{moved} into the struct (unless
it has \texttt{Copy} type). After construction, any non-\texttt{Copy} variables
used in $e_i$ are marked as moved.

\paragraph{Field Access}

\[
\frac{
  \Gamma ; \Sigma \vdash e : S \quad \mathit{fields}(S)(f) = \tau
}{
  \Gamma ; \Sigma \vdash (.~e~f) : \tau
}
\quad \textsc{T-Field}
\]

Accessing a field of a struct returns a value of the field's type. If $e$ is an
owned struct, the entire struct is consumed (moved). If $e$ is a reference
$\&S$ or $\&_{mut}S$, the field access borrows from the reference.

\paragraph{Immutable Borrowing}

\[
\frac{
  \Gamma(x) = \tau \quad \Sigma(x) = owned
}{
  \Gamma ; \Sigma[x \mapsto borrowed(1)] \vdash (\mathtt{borrow}\ x) : \&\tau
}
\quad \textsc{T-Borrow}
\]

\[
\frac{
  \Gamma(x) = \tau \quad \Sigma(x) = borrowed(n)
}{
  \Gamma ; \Sigma[x \mapsto borrowed(n+1)] \vdash (\mathtt{borrow}\ x) : \&\tau
}
\quad \textsc{T-Borrow-Shared}
\]

Borrowing creates an immutable reference without transferring ownership. Multiple
immutable borrows can coexist (incrementing the borrow count). While borrowed,
the original variable cannot be moved or mutably borrowed, preventing
use-after-free and data races.

\paragraph{Mutable Borrowing}

\[
\frac{
  \Gamma(x) = \tau \quad \Sigma(x) = owned
}{
  \Gamma ; \Sigma[x \mapsto borrowed_{mut}] \vdash (\mathtt{borrow\text{-}mut}\ x) : \&_{mut}\tau
}
\quad \textsc{T-Borrow-Mut}
\]

A mutable borrow is \emph{exclusive}: no other borrows (mutable or immutable) can
exist simultaneously. This guarantees that mutations through the reference cannot
create data races or iterator invalidation.

\paragraph{Function Application (Move Semantics)}

\[
\frac{
  \Gamma ; \Sigma \vdash f : (\tau_1 \dots \tau_n \rightarrow \tau_r) \quad
  \Gamma ; \Sigma_0 \vdash e_i : \tau_i \quad
  \Sigma_{i+1} = \Sigma_i[e_i \mapsto moved]
}{
  \Gamma ; \Sigma_n \vdash (f\ e_1 \dots e_n) : \tau_r
}
\quad \textsc{T-Call}
\]

When calling a function with owned parameters, arguments are moved into the
function. Each argument evaluation potentially changes the ownership state
(sequentially threaded through $\Sigma_i$). After the call, moved variables
are inaccessible.

\paragraph{Function Application (Borrow Semantics)}

\[
\frac{
  \Gamma ; \Sigma \vdash f : (\&\tau \rightarrow \tau_r) \quad
  \Gamma ; \Sigma \vdash e : \&\tau
}{
  \Gamma ; \Sigma \vdash (f\ e) : \tau_r
}
\quad \textsc{T-Call-Borrow}
\]

Functions accepting references do not consume their arguments. The caller retains
ownership, and the ownership state $\Sigma$ is unchanged (aside from any borrow
state changes from creating the reference).


\subsection{Operational Semantics}

Ferrite uses call-by-value evaluation with explicit ownership transitions. The
operational semantics model how expressions reduce to values and how ownership
states evolve during execution.

\paragraph{Variable Evaluation (Move)}

\[
\frac{
  \Sigma(x) = owned
}{
  \langle x, \Sigma \rangle \rightarrow \langle v, \Sigma[x \mapsto moved] \rangle
}
\quad \textsc{E-Var-Move}
\]

Reading an owned variable \emph{moves} its value, leaving the variable in a moved
state. This implements affine types: each owned value is used at most once.

\paragraph{Variable Evaluation (Copy)}

\[
\frac{
  \Sigma(x) = owned \quad \Gamma(x) = \tau \quad \tau \in \mathit{Copy}
}{
  \langle x, \Sigma \rangle \rightarrow \langle v, \Sigma \rangle
}
\quad \textsc{E-Var-Copy}
\]

For \texttt{Copy} types, reading a variable produces a copy of the value without
affecting the ownership state.

\paragraph{Borrow Evaluation}

\[
\langle (\mathtt{borrow}\ x), \Sigma \rangle
\rightarrow
\langle \&x, \Sigma[x \mapsto borrowed(n+1)] \rangle
\quad \textsc{E-Borrow}
\]

Creating a borrow produces a reference value and increments the borrow count.
The variable remains accessible for further immutable borrows.

\paragraph{Let Evaluation (Step)}

\[
\frac{
  \langle e_1, \Sigma \rangle \rightarrow \langle e_1', \Sigma' \rangle
}{
  \langle (\mathtt{let}\ ((x\ e_1))\ e_2), \Sigma \rangle
  \rightarrow
  \langle (\mathtt{let}\ ((x\ e_1'))\ e_2), \Sigma' \rangle
}
\quad \textsc{E-Let-Step}
\]

If the bound expression can take a step, the let-expression steps by evaluating
the binding expression.

\paragraph{Let Evaluation (Bind)}

\[
\langle (\mathtt{let}\ ((x\ v))\ e), \Sigma \rangle
\rightarrow
\langle e[x \mapsto v], \Sigma[x \mapsto owned] \rangle
\quad \textsc{E-Let-Bind}
\]

Once the bound expression is a value, substitute it into the body and mark the
variable as owned in the ownership environment.

\paragraph{Struct Construction}

\[
\frac{
  \forall i : \langle e_i, \Sigma_{i-1} \rangle \rightarrow \langle v_i, \Sigma_i \rangle
}{
  \langle (S\ e_1 \dots e_n), \Sigma_0 \rangle
  \rightarrow
  \langle S(v_1,\dots,v_n), \Sigma_n \rangle
}
\quad \textsc{E-Struct}
\]

Struct construction evaluates each field expression left-to-right, threading
ownership state changes through the sequence.


\subsection{Scope Exit}

When execution exits a scope (e.g., at the end of a let-body), any active borrows
are released:

\[
\Sigma(x) = borrowed(n) \Rightarrow
\Sigma'(x) =
\begin{cases}
owned & \text{if } n = 1 \\
borrowed(n-1) & \text{if } n > 1
\end{cases}
\]

This models the automatic release of borrows when their lexical scope ends.
If the variable was moved or is still borrowed elsewhere, it retains that state.


\subsection{Type Soundness}

Type soundness is the fundamental safety property: well-typed programs do not
get stuck in undefined states. We prove this through the standard progress and
preservation approach.

\subsubsection{Lemma 1 (Canonical Forms)}

If $\Gamma ; \Sigma \vdash v : \tau$ and $v$ is a value, then:

\begin{itemize}
  \item If $\tau = i32$, then $v = n$ for some integer $n$
  \item If $\tau = bool$, then $v \in \{true, false\}$
  \item If $\tau = S$, then $v = S(v_1,\dots,v_n)$ where each $v_i$ is a value
  \item If $\tau = \&\tau'$, then $v = \&x$ for some variable $x$ with $\Gamma(x) = \tau'$
\end{itemize}

\paragraph{Proof Sketch}
By case analysis on the typing derivation. Each typing rule for values ensures
the value has the expected form. For example, \textsc{T-Borrow} is the only rule
that produces type $\&\tau$, and it requires the expression to be
$(\mathtt{borrow}\ x)$, which evaluates to $\&x$.


\subsubsection{Theorem 1 (Progress)}

If $\emptyset ; \Sigma \vdash e : \tau$, then either:
\begin{enumerate}
  \item $e$ is a value, or
  \item there exist $e', \Sigma'$ such that
    $\langle e, \Sigma \rangle \rightarrow \langle e', \Sigma' \rangle$
\end{enumerate}

\paragraph{Proof Sketch}
By induction on the typing derivation of $e$.

\textbf{Base cases:}
\begin{itemize}
  \item Literals ($n$, $true$, $false$) are values
  \item Variables: By \textsc{T-Var}, $\Sigma(x) = owned$, so \textsc{E-Var-Move}
    applies
\end{itemize}

\textbf{Inductive cases:}
\begin{itemize}
  \item Let-binding: By induction, $e_1$ is either a value (so \textsc{E-Let-Bind}
    applies) or steps (so \textsc{E-Let-Step} applies)
  \item Borrow: By typing rule, variable is owned or borrowed, so
    \textsc{E-Borrow} applies
  \item Function call: By induction, arguments step or are values; when all are
    values, the function body steps
\end{itemize}

Well-typed expressions cannot get stuck because ownership guarantees in $\Sigma$
align exactly with the preconditions of evaluation rules.


\subsubsection{Theorem 2 (Preservation)}

If $\Gamma ; \Sigma \vdash e : \tau$ and
$\langle e, \Sigma \rangle \rightarrow \langle e', \Sigma' \rangle$,
then $\Gamma ; \Sigma' \vdash e' : \tau$.

\paragraph{Proof Sketch}
By induction on the evaluation step.

\textbf{Key cases:}
\begin{itemize}
  \item \textsc{E-Var-Move}: Type unchanged ($\tau$), ownership updated to
    $moved$. Since the variable is no longer accessible, this preserves
    well-typedness.
  \item \textsc{E-Borrow}: Type becomes $\&\tau$, ownership becomes
    $borrowed(n+1)$. The typing rule \textsc{T-Borrow} explicitly allows this
    transition.
  \item \textsc{E-Let-Bind}: Substitution lemma ensures $e[x \mapsto v]$ has the
    same type as $e$ under the extended environment. The ownership state correctly
    reflects that $x$ is now owned.
  \item \textsc{E-Struct}: Each field steps preserving its type, so the final
    struct value has the correct type.
\end{itemize}

The critical insight is that ownership transitions in $\Sigma$ mirror the
constraints in the typing rules, so type judgements are preserved across
evaluation steps.


\subsubsection{Corollary (Type Soundness)}

Well-typed Ferrite programs do not exhibit:
\begin{itemize}
  \item \textbf{Use-after-move}: Variables marked as $moved$ cannot be used
    (enforced by \textsc{T-Var})
  \item \textbf{Dangling references}: References cannot outlive their referents
    (enforced by lexical scoping and borrow rules)
  \item \textbf{Data races}: Mutable borrows are exclusive; immutable borrows
    prevent mutation
  \item \textbf{Invalid field access}: Field access requires the struct type to
    define that field
\end{itemize}

\paragraph{Proof}
These properties follow directly from Progress (no well-typed program gets stuck)
and Preservation (types and ownership are maintained during execution). Any
violation would require either a stuck state (contradicting Progress) or a type
mismatch after a step (contradicting Preservation).


\subsection{Traits and Polymorphism}

Ferrite supports trait-based polymorphism, enabling generic programming while
maintaining type safety.

\paragraph{Trait Definitions}
A trait $T$ declares a set of method signatures:
\[
\mathit{trait}\ T = \{m_i : \tau_i\}
\]

\paragraph{Trait Implementation}
A type $S$ can implement trait $T$ by providing definitions for all methods:
\[
\mathit{impl}\ T\ \mathit{for}\ S
\]

\paragraph{Trait Bounds}
Function types can be constrained by trait bounds:
\[
\forall \alpha : T . (\alpha \rightarrow \tau)
\]
This reads: ``for any type $\alpha$ implementing trait $T$, a function from
$\alpha$ to $\tau$.''

\paragraph{Built-in Traits}
\begin{itemize}
  \item $\mathit{Copy}$: Types whose values can be duplicated bitwise
  \item $\mathit{Clone}$: Types that support explicit cloning
  \item $\mathit{Debug}$: Types that can be formatted for debugging
\end{itemize}


\subsection{Conclusion}

Ferrite's safety guarantees arise from the synergy of three key design elements:

\begin{enumerate}
  \item \textbf{Affine types with Copy}: Values are used at most once unless
    explicitly marked as copyable, preventing use-after-free at compile time.
  
  \item \textbf{Explicit borrow tracking}: References are first-class, with
    ownership states enforcing aliasing XOR mutability (you can have many
    immutable borrows OR one mutable borrow, but not both).
  
  \item \textbf{Operational ownership transitions}: The operational semantics
    directly model ownership transfer and borrow creation, making the runtime
    behavior of the ownership system explicit and verifiable.
\end{enumerate}

The language is both \textbf{implementable} (the type system is decidable and
compiles to efficient C code) and \textbf{formally sound} (proven memory-safe via
Progress and Preservation). This makes Ferrite suitable for systems programming
where safety and performance are both critical requirements.

Future extensions could include:
\begin{itemize}
  \item Lifetime parameters for more flexible borrowing patterns
  \item Higher-kinded types for advanced trait abstractions
  \item Effect systems for tracking computational effects (I/O, exceptions)
  \item Dependent types for richer compile-time guarantees
\end{itemize}


\end{document}
